% 61e Olympiade Internationale de Mathématiques — Saint-Pétersbourg, Russie.
% Énoncés officiels en français (imo-official.org). Chaque problème vaut 7 points.

\begin{center}
  \emph{Premier jour --- lundi 21 septembre 2020 --- Durée : 4 heures 30}
\end{center}

\subsection*{Problème 1}
On considère un quadrilatère convexe $ABCD$. Un point $P$ est situé à
l'intérieur de $ABCD$. On suppose que les égalités de rapports ci-dessous
sont vérifiées :
\[
\widehat{PAD} : \widehat{PBA} : \widehat{DPA}
= 1 : 2 : 3 =
\widehat{CBP} : \widehat{BAP} : \widehat{BPC}.
\]
Montrer que les trois droites suivantes se rencontrent en un point : la
bissectrice intérieure de l'angle $\widehat{ADP}$, la bissectrice intérieure
de l'angle $\widehat{PCB}$ et la médiatrice du segment $[AB]$.

\subsection*{Problème 2}
Soit $a$, $b$, $c$, $d$ des nombres réels tels que $a \geqslant b \geqslant c
\geqslant d > 0$ et $a + b + c + d = 1$. Montrer que
\[
(a + 2b + 3c + 4d)\, a^{a}\, b^{b}\, c^{c}\, d^{d} < 1.
\]

\subsection*{Problème 3}
On se donne $4n$ cailloux de poids $1, 2, 3, \ldots, 4n$. Chaque caillou est
coloré en une couleur parmi $n$ couleurs possibles, et il y a quatre cailloux
de chaque couleur. Montrer que l'on peut répartir les cailloux en deux tas de
sorte que les deux conditions suivantes soient vérifiées :
\begin{itemize}
  \item Les poids totaux des deux tas sont égaux.
  \item Chaque tas contient deux cailloux de chaque couleur.
\end{itemize}

\bigskip
\begin{center}
  \emph{Second jour --- mardi 22 septembre 2020 --- Durée : 4 heures 30}
\end{center}

\subsection*{Problème 4}
Soit $n > 1$ un entier. Il y a $n^2$ stations sur le versant d'une montagne,
toutes à des altitudes différentes. Chacune des deux compagnies de
téléphériques, $A$ et $B$, gère $k$ téléphériques ; chaque téléphérique
permet de se déplacer d'une des stations vers une station plus élevée (sans
arrêt intermédiaire). Les $k$ téléphériques de $A$ ont $k$ points de départ
différents et $k$ points d'arrivée différents, et un téléphérique qui a un
point de départ plus élevé a aussi un point d'arrivée plus élevé. Les mêmes
conditions sont satisfaites pour $B$. On dit que deux stations sont
\emph{reliées} par une compagnie s'il est possible de partir de la station la
plus basse et d'atteindre la plus élevée en utilisant un ou plusieurs
téléphériques de cette compagnie (aucun autre mouvement entre les stations
n'est autorisé).

Déterminer le plus petit entier strictement positif $k$ qui garantisse qu'il
existe deux stations reliées par chacune des deux compagnies.

\subsection*{Problème 5}
On se donne un paquet de $n > 1$ cartes. Un entier strictement positif est
écrit sur chaque carte. Le paquet a la propriété que la moyenne arithmétique
des nombres écrits sur chaque paire de cartes est aussi égale à la moyenne
géométrique des nombres écrits sur une certaine collection d'une ou plusieurs
cartes.

Pour quels $n$ cela implique-t-il que les nombres écrits sur les cartes sont
tous égaux ?

\subsection*{Problème 6}
Montrer qu'il existe une constante strictement positive $c$ telle que
l'assertion suivante soit vraie :

Considérons un entier $n > 1$, et un ensemble $S$ de $n$ points du plan tel
que la distance entre deux points distincts dans $S$ soit toujours au moins
égale à $1$. Alors il existe une droite $\ell$ qui sépare $S$, de sorte que la
distance de n'importe quel point de $S$ à $\ell$ vaut au moins $c\, n^{-1/3}$.

(On dit qu'une droite $\ell$ \emph{sépare} un ensemble de points $S$ lorsque
$\ell$ rencontre au moins un segment joignant deux points de $S$.)

\medskip
\emph{Remarque.} Des résultats plus faibles avec $c\, n^{-\alpha}$ à la place
de $c\, n^{-1/3}$ pourront recevoir des points selon la valeur de la constante
$\alpha > 1/3$.
